\documentclass[11pt]{article}
\usepackage[T1]{fontenc}
\usepackage{lmodern,microtype,amsmath,amssymb,amsthm,mathtools,mathrsfs,tikz-cd}
\usepackage[margin=1in]{geometry}
\usepackage{enumitem,booktabs,longtable,xurl,hyperref}
\hypersetup{colorlinks=true,linkcolor=blue,citecolor=blue,urlcolor=blue,
 pdftitle={Domains, finite proofs, and essential incompleteness},pdfauthor={Florian Lengyel}}
\newtheorem{theorem}{Theorem}[section]
\newtheorem{proposition}[theorem]{Proposition}
\newtheorem{lemma}[theorem]{Lemma}
\newtheorem{corollary}[theorem]{Corollary}
\theoremstyle{definition}
\newtheorem{definition}[theorem]{Definition}
\newtheorem{example}[theorem]{Example}
\theoremstyle{remark}
\newtheorem{remark}[theorem]{Remark}
\newtheorem{problem}[theorem]{Problem}
\newcommand{\R}{\mathscr R}
\newcommand{\D}{\mathsf D}
\newcommand{\Dom}{\operatorname{Dom}}
\newcommand{\ran}{\operatorname{ran}}
\newcommand{\Tot}{\operatorname{Tot}}
\newcommand{\Pos}{\mathbf{Pos}}
\newcommand{\Set}{\mathbf{Set}}
\newcommand{\Prof}{\mathbf{Prof}}
\newcommand{\CE}{\operatorname{CE}}
\newcommand{\N}{\mathbb N}
\newcommand{\op}{\mathrm{op}}
\newcommand{\id}{\operatorname{id}}
\newcommand{\Cons}{\mathsf{Cons}}
\newcommand{\Th}{\mathsf{Th}}
\newcommand{\BDLat}{\mathbf{BDLat}}
\title{Domains, finite proofs,\\and essential incompleteness}
\author{Florian Lengyel}
\date{29 October 2026}
\begin{document}
\maketitle
\begin{abstract}
Restriction domains and ranges define profunctors for total and partial
correctness. For partial recursive maps, the partial-correctness profunctor
has representing computational domains exactly when the program domain is
recursive. A Turing morphism, together with binary case distinction,
constructs an inseparable pair of domains. Typed inference rules determine
polynomial functors. For recursive presentations with at most two premises
per rule, finite proofs form their initial algebras in the total category
obtained by splitting numerical restriction idempotents.
Taking the range of a proof-conclusion map gives a functorial theorem
object. Computation traces provide positive and negative proofs in an
explicit classical propositional presentation, and every consistent
computably enumerable extension of that presentation is incomplete.
The corresponding existence problem for logical presentations in general
recursion categories is formulated separately.
\end{abstract}
\tableofcontents

\section{Computation, assertions, and deduction}
Essential incompleteness concerns consistent effective extensions of a
deductive system. A categorical account must therefore connect computational
structure with assertions and deductions, and explain how consistency and
completeness impose incompatible conditions on computational domains.
Di Paola and Heller formulate this problem through categorical analogues of
recursion-theoretic inseparability and creativity, together with an intrinsic
connection to algebraic formal systems \cite[\S1]{DPH}.

A restriction domain describes where a partial computation is defined.
A Hoare relation describes which preconditions and postconditions that
computation satisfies. A polynomial inference functor describes how
premise proofs form a proof of a conclusion. These constructions have
different categorical types. The domain construction is contravariant;
the Hoare construction is a covariant pseudofunctor to profunctors; and
the finite-proof construction is an initial algebra in a slice category.
Ranges then associate computational theorem objects to deductive
presentations.

The argument for essential incompleteness has three steps. Turing
universality and binary case distinction produce inseparable computational
domains. Maps from finite computation evidence to finite proofs transfer
these domains to provability and refutability. Consistency and completeness
would turn those two domains into a complementary separator. The numerical
construction supplies all three steps in a category of computational
objects. For a general recursion category, the existence of the requisite
logical presentations remains a further question.

Predicate-transformer semantics associates programs with transformations
of assertions \cite{Dijkstra,DijkstraScholten}. The Hoare constructions
below use the local direct-image and inverse-image adjunction of range
restriction categories \cite[Proposition~8.2]{CockettGarner}, and have an
antecedent in the categorical framework of Arthan, Martin, Mathiesen, and
Oliva \cite{ArthanEtAl}. Heller's stable unions supply the domain
operations used in separation \cite[\S6]{Heller}. Polynomial bicomodules
organize the premise and conclusion types of rules
\cite[\S8.3]{NS2024}; their relation to ordinary profunctors is described
through right adjoints between functor categories. The final arithmetic
corollary assumes two-sided numeralwise definability of recursive sets.

Composition is written $gf=g\circ f$, with $f$ performed first. Posets are
regarded as categories with one arrow $p\to q$ precisely when $p\le q$.
All categories used to index coends are small, or are taken in a fixed
universe in which the indicated coends exist. Effectivity in the numerical
sections refers to ordinary computability on natural-number codes.


\section{Restriction domains and direct images}\label{hp:domains}

A \emph{restriction
category} is a category $\R$ equipped with an operation assigning to each
$f:X\to Y$ an endomorphism $\bar f:X\to X$ such that
\begin{equation}\label{hp:restriction-axioms}
 f\bar f=f,\qquad
 \bar f\bar g=\bar g\bar f,\qquad
 \overline{g\bar f}=\bar g\bar f,\qquad
 \bar h f=f\overline{hf}.
\end{equation}
In the second and third identities $f$ and $g$ have the same source; in
the fourth, $f:X\to Y$ and $h:Y\to Z$. An arrow $f:X\to Y$ is
\emph{total} if $\bar f=1_X$. Total arrows form a subcategory
$\Tot(\R)$ with the same objects as $\R$. We assume local smallness and
zero arrows satisfying $\overline{0_{X,Y}}=0_{X,X}$. The restriction
structure is the part of a recursion category used in this section
\cite{CockettLack}.

A \emph{restriction idempotent} on $X$ is an endomorphism $p:X\to X$
with $\bar p=p$. It expresses a domain on $X$. Let $\Dom(X)$ be the set
of these idempotents, ordered by
\[
 p\le q\quad\Longleftrightarrow\quad qp=p.
\]
The meet is $p\wedge q=pq$, the greatest element is $1_X$, and the
least element is $0_X$. Regard this poset as a category: its hom-set
$\Dom(X)(p,q)$ is a singleton if $p\le q$ and is empty otherwise.

Define the contravariant functor
\begin{equation}\label{hp:domain-functor}
 \D:\R^{\op}\longrightarrow\Pos.
\end{equation}
On objects put $\D(X)=\Dom(X)$. For $f:X\to Y$, the monotone map
$\D(f^{\op}):\D(Y)\to\D(X)$ is denoted $f^*$ and is given by
\[
 f^*(q)=\overline{qf}.
\]
The restriction identities imply
\begin{equation}\label{hp:pullback-laws}
 (gf)^*=f^*g^*,\qquad 1_X^*=1_{\D(X)},\qquad
 f^*(pq)=f^*(p)f^*(q),\qquad
 f^*(1_Y)=\bar f,\qquad f^*(0_Y)=0_X.
\end{equation}
These are the identity, composition, and meet laws of the fundamental
domain construction; see \cite[Proposition~3.1]{CockettGarner}.
In particular, inverse image preserves the greatest element precisely
for total arrows.

Suppose that $\R$ has a \emph{range operator}. Its value on $f:X\to Y$
is a restriction idempotent $\ran(f):Y\to Y$ satisfying
\begin{equation}\label{hp:range-axioms}
 \ran(f)f=f,\qquad
 \ran(ef)=e\ran(f),\qquad
 \ran(g\ran(f))=\ran(gf),
\end{equation}
where $e$ is a restriction idempotent on $Y$ and $g:Y\to Z$.
These are the range-restriction axioms, with the assertion that
$\ran(f)$ is a restriction idempotent included in the definition
\cite[Section~8]{CockettGarner}. They give the following universal
property among domains:
\begin{equation}\label{hp:range-least}
 \ran(f)\le e\quad\Longleftrightarrow\quad ef=f.
\end{equation}
Indeed, the forward implication follows from $\ran(f)f=f$, and the
reverse implication follows from $\ran(f)=\ran(ef)=e\ran(f)$.

Define the covariant functor
\[
 \D_!:\R\longrightarrow\Pos
\]
by $\D_!(X)=\D(X)$ and
\[
 \D_!(f)=f_!: \D(X)\longrightarrow\D(Y),\qquad
 f_!(p)=\ran(fp).
\]
If $p\le p'$, every domain containing the range of $fp'$ also contains
the range of $fp$, so \eqref{hp:range-least} proves monotonicity.
Since $\ran(p)=p$ for a restriction idempotent, and by the last range
axiom,
\begin{equation}\label{hp:image-functor}
 (1_X)_!=1_{\D(X)},\qquad (gf)_!=g_!f_!.
\end{equation}

\begin{proposition}[Local adjunction]\label{hp:local-adjunction}
Let $f:X\to Y$ be an arrow in a range restriction category. Restricting
the domain of $f_!$ and the codomain of $f^*$ gives an adjunction
\[
 f_!:\D(X)\mathbin{\downarrow}\bar f
 \ \rightleftarrows\ \D(Y):f^*,\qquad f_!\dashv f^*,
\]
where $\D(X)\mathbin{\downarrow}\bar f=\{p:p\le\bar f\}$.
For arbitrary $p\in\D(X)$ and $q\in\D(Y)$,
\begin{equation}\label{hp:partial-adjunction}
 f_!(p)\le q\quad\Longleftrightarrow\quad
 p\bar f\le f^*(q).
\end{equation}
Moreover, direct image satisfies Frobenius reciprocity:
\begin{equation}\label{hp:frobenius}
 f_!\bigl(p\wedge f^*(q)\bigr)=f_!(p)\wedge q.
\end{equation}
\end{proposition}
\begin{proof}
First suppose $p\le\bar f$. By \eqref{hp:range-least},
$f_!(p)\le q$ is equivalent to $qfp=fp$. If $p\le\overline{qf}$,
the last restriction identity gives $qfp=f\overline{qf}p=fp$.
Conversely, taking restrictions of $qfp=fp$ gives
$\overline{qf}p=p$, since $p\le\bar f$. This proves the adjunction.
For arbitrary $p$, use $fp=f\bar f p$ to obtain
\eqref{hp:partial-adjunction}. Finally,
$fp\overline{qf}=qfp$, so
\[
 f_!\bigl(p\wedge f^*(q)\bigr)
 =\ran(qfp)=q\ran(fp).
\]
This is also the local adjunction of
\cite[Proposition~8.2]{CockettGarner}.
\end{proof}

\section{Hoare assertions as profunctors}\label{hp:profunctors}

For small categories $A,B$, a \emph{profunctor} $P:A\nrightarrow B$
means a functor
\[
 P:A^{\op}\times B\longrightarrow\Set.
\]
For $P:A\nrightarrow B$ and $Q:B\nrightarrow C$, their composite is
\[
 (Q\odot P)(a,c)=\int^{b\in B}P(a,b)\times Q(b,c).
\]
Here the coend is the quotient of the disjoint union of the displayed
products by the identifications induced by arrows of $B$: for
$u:b\to b'$, $x\in P(a,b)$, and $y\in Q(b',c)$,
\[
 [P(1,u)x,y]=[x,Q(u^{\op},1)y].
\]
The identity profunctor on $A$ is its hom-functor $A(-,-)$.
Categories, profunctors, and natural transformations form the bicategory
$\Prof$; its associativity constraints are the canonical reassociations
of these coends. Each $\D(X)$ above is small because $\R$ is locally
small.

For an arrow $f:X\to Y$, define the functors
\begin{equation}\label{hp:hoare-types}
 H_f,K_f:\D(X)^{\op}\times\D(Y)\longrightarrow\Set
\end{equation}
by
\begin{equation}\label{hp:hoare-definitions}
 H_f(p,q)=\D(X)(p,f^*q),\qquad
 K_f(p,q)=\D(Y)(f_!p,q).
\end{equation}
Their values are empty sets or singletons. On morphisms, both functors
use strengthening of preconditions and weakening of postconditions:
if $p'\le p$ and $q\le q'$, the unique maps
\[
 H_f(p,q)\longrightarrow H_f(p',q'),\qquad
 K_f(p,q)\longrightarrow K_f(p',q')
\]
exist by monotonicity. This specifies their functorial action.

In the category of sets and partial functions, $H_f(p,q)$ is inhabited
exactly when $f$ terminates on every input in $p$ and every resulting
output lies in $q$. The value $K_f(p,q)$ is inhabited exactly when every
terminating computation from $p$ has its output in $q$. These are,
respectively, total and partial correctness. The definitions in
\eqref{hp:hoare-definitions} express the same distinction using only
domains, inverse image, and range.

\begin{proposition}\label{hp:hoare-composition}
The assignments $X\mapsto\D(X)$, $f\mapsto H_f$ and
$X\mapsto\D(X)$, $f\mapsto K_f$ define covariant pseudofunctors
\[
 H,K:\R\longrightarrow\Prof,
\]
where $\R$ is regarded as a bicategory with only identity $2$-cells.
Their composition constraints are
\[
 H_g\odot H_f\cong H_{gf},\qquad
 K_g\odot K_f\cong K_{gf}.
\]
\end{proposition}
\begin{proof}
For $f:X\to Y$, $g:Y\to Z$, $p\in\D(X)$, and $r\in\D(Z)$,
the co-Yoneda identity gives
\begin{align*}
 (H_g\odot H_f)(p,r)
 &=\int^{q\in\D(Y)}\D(X)(p,f^*q)\times\D(Y)(q,g^*r)\\
 &\cong\D(X)(p,f^*g^*r)=H_{gf}(p,r),\\
 (K_g\odot K_f)(p,r)
 &=\int^{q\in\D(Y)}\D(Y)(f_!p,q)\times\D(Z)(g_!q,r)\\
 &\cong\D(Z)(g_!f_!p,r)=K_{gf}(p,r).
\end{align*}
For completeness, the first isomorphism sends a pair of inequalities
$p\le f^*q$, $q\le g^*r$ to their composite
$p\le f^*g^*r$. Its inverse uses $q=g^*r$; every other factorization
is identified with this one by the coend relation. The second
isomorphism composes $f_!p\le q$ with $g_!q\le r$ after applying
$g_!$ to the first inequality; its inverse uses $q=f_!p$.
Identity arrows give the hom-profunctors. The associativity and unit
identities hold because both routes compose the same inequalities.
\end{proof}

The covariance of these pseudofunctors is compatible with the
contravariance of $\D$: the right-hand argument $f^*q$ of the hom-set in
$H_f$ converts inverse image into a profunctor from $\D(X)$ to
$\D(Y)$.

\begin{proposition}\label{hp:hoare-comparison}
For $f:X\to Y$ and domains $p,q$,
\begin{align}
 K_f(p,q)\ne\varnothing
 &\quad\Longleftrightarrow\quad p\bar f\le f^*q,
 \label{hp:k-equivalence}\\
 H_f(p,q)\ne\varnothing
 &\quad\Longleftrightarrow\quad
 p\le\bar f\ \text{ and }\ K_f(p,q)\ne\varnothing.
 \label{hp:h-equivalence}
\end{align}
There is a canonical natural transformation $H_f\Rightarrow K_f$.
It is an isomorphism if and only if $f$ is total.
\end{proposition}
\begin{proof}
Equation~\eqref{hp:k-equivalence} is
\eqref{hp:partial-adjunction}. Since $f^*q\le\bar f$,
the condition $p\le f^*q$ is equivalent to the two conditions in
\eqref{hp:h-equivalence}. This gives the natural transformation.
If $f$ is total, the two conditions coincide for all $p,q$.
Conversely, $K_f(1_X,1_Y)$ is always inhabited, whereas
$H_f(1_X,1_Y)$ is inhabited exactly when $1_X\le\bar f$.
\end{proof}

For the zero arrow $0_{X,Y}$, $H_0(p,q)$ is inhabited exactly when
$p=0_X$, whereas $K_0(p,q)$ is always inhabited. Thus the distinction
already appears for a nowhere-defined program.

The restriction order on parallel arrows is
$f\le g$ if $f=g\bar f$. It induces opposite orders on the two
profunctor assignments:
\begin{equation}\label{hp:program-order}
 f\le g\quad\Longrightarrow\quad
 H_f\Rightarrow H_g\quad\text{and}\quad K_g\Rightarrow K_f.
\end{equation}
Indeed, $f^*q=\bar f g^*q\le g^*q$, and
$f_!p=\ran(g\bar fp)\le\ran(gp)=g_!p$. Each displayed natural
transformation is the unique map between the corresponding subterminal
values. Thus extension of a partial program enlarges its total-correctness
relation and restricts its partial-correctness relation.

\begin{proposition}[Detection of domains]\label{hp:detect-domains}
For restriction idempotents $e,d$ on $X$,
\[
 H_e(p,1_X)\ne\varnothing\quad\Longleftrightarrow\quad p\le e,
 \qquad
 H_e\Rightarrow H_d\quad\Longleftrightarrow\quad e\le d.
\]
Consequently $H_e\cong H_d$ implies $e=d$, and
\[
 ed=0_X\quad\Longleftrightarrow\quad H_d\odot H_e\cong H_{0_X}.
\]
\end{proposition}
\begin{proof}
The first equivalence follows from $e^*(1_X)=e$.
If $H_e\Rightarrow H_d$, evaluate it at $(e,1_X)$ to obtain $e\le d$.
The converse follows from \eqref{hp:program-order}. The final statement
uses $H_d\odot H_e\cong H_{de}$ and the fact that restriction
idempotents commute.
\end{proof}

For partial recursive maps on $\N$, singleton domains distinguish arrows: $H_f(\{x\},\{y\})$ is inhabited exactly when
$f(x)=y$. Hence $H_f\cong H_g$ implies $f=g$ in that category.
More precisely, $H_f\Rightarrow H_g$ holds exactly when $f\le g$.

\section{Weakest liberal preconditions and representability}
\label{hp:representability}

For each $q\in\D(Y)$, partial correctness determines the presheaf
\[
 K_f(-,q):\D(X)^{\op}\longrightarrow\Set.
\]
A \emph{weakest liberal precondition} for $q$ is a representing domain
$w\in\D(X)$, that is, a natural isomorphism
\begin{equation}\label{hp:wlp-representing}
 K_f(p,q)\cong\D(X)(p,w)
 \qquad(p\in\D(X)).
\end{equation}
Equivalently, $w$ is the greatest domain satisfying $f_!w\le q$.
If such a domain exists for every $q$, it defines a monotone map
$\operatorname{wlp}_f:\D(Y)\to\D(X)$ and an adjunction
\[
 f_!:\D(X)\rightleftarrows\D(Y):\operatorname{wlp}_f.
\]
The domain of this left adjoint is all of $\D(X)$; compare the local
adjunction in Proposition~\ref{hp:local-adjunction}.

Suppose that $\D(X)$ has finite joins distributing over meets, and that
$\bar f$ has a complement $c$, meaning $c\bar f=0_X$ and
$c\vee\bar f=1_X$. Then every weakest liberal precondition exists and
\begin{equation}\label{hp:complemented-wlp}
 \operatorname{wlp}_f(q)=c\vee f^*q.
\end{equation}
Indeed, distributivity gives
$p\bar f\le f^*q$ if and only if $p\le c\vee f^*q$; apply
\eqref{hp:k-equivalence}. Stable finite unions of domains supply the
distributivity used here.

Existence of a weakest liberal precondition for $0_Y$ alone means that
$\bar f$ has a \emph{pseudocomplement}: a greatest domain disjoint from
$\bar f$. It does not imply the complement equation
$\bar f\vee\operatorname{wlp}_f(0_Y)=1_X$ in an arbitrary domain
lattice. For example, in the chain $0<a<1$, the pseudocomplement of
$a$ is $0$, and $a\vee0\ne1$. This is a domain lattice: take the
one-object category whose arrows are $0,a,1$, whose composition is meet,
and whose restriction and range operators both fix every arrow.

\subsection{The numerical obstruction}

A subset of $\N$ is \emph{recursive} if membership is decidable, and
\emph{computably enumerable} (c.e.) if it is the domain of a partial
recursive function. Let $\CE(\N)$ be the poset of c.e. subsets of
$\N$, ordered by inclusion. These are precisely the restriction
domains on $\N$ in the category of partial recursive maps: a c.e.
set $P$ determines the partial identity with domain $P$, and every
restriction idempotent is such a partial identity. For a partial
recursive $f:\N\rightharpoonup\N$, both $f(P)$ and
$\{n:f(n)\text{ is defined and belongs to }Q\}$ are c.e. when
$P,Q$ are c.e. They give the maps $f_!$ and $f^*$ on this poset.

\begin{theorem}\label{hp:numerical-wlp}
For a partial recursive function $f:\N\rightharpoonup\N$, the map
\[
 f_!:\CE(\N)\longrightarrow\CE(\N),\qquad P\longmapsto f(P),
\]
has a right adjoint if and only if $\operatorname{dom}f$ is recursive.
In fact, representability of $K_f(-,\varnothing)$ already implies that
$\operatorname{dom}f$ is recursive.
\end{theorem}
\begin{proof}
Suppose that a c.e. set $W$ represents $K_f(-,\varnothing)$. Then
\[
 P\subseteq W\quad\Longleftrightarrow\quad
 P\cap\operatorname{dom}f=\varnothing
 \qquad(P\in\CE(\N)).
\]
Taking $P=W$ shows that $W$ is disjoint from $\operatorname{dom}f$.
If $n\notin\operatorname{dom}f$, taking the recursive singleton
$P=\{n\}$ gives $n\in W$. Thus
$W=\N\setminus\operatorname{dom}f$. The domain of a partial recursive
function is c.e.; since its complement is also c.e., dovetailing the two
enumerations decides membership in $\operatorname{dom}f$.

Conversely, if $\operatorname{dom}f$ is recursive, then
\[
 W_Q=(\N\setminus\operatorname{dom}f)\cup f^*Q
\]
is c.e. for every c.e. $Q$. For every c.e. $P$,
$f(P)\subseteq Q$ holds exactly when $P\subseteq W_Q$.
Therefore $Q\mapsto W_Q$ is the required right adjoint.
\end{proof}

Thus a partial-correctness presheaf can exist without a representing
computational domain. At the false postcondition its values describe
all c.e. domains of nontermination, even when their union is not c.e.:
\[
 K_f(P,\varnothing)\ne\varnothing
 \quad\Longleftrightarrow\quad P\cap\operatorname{dom}f=\varnothing.
\]
The use of singleton domains in the proof explains why the numerical
necessity statement is stronger than the abstract pseudocomplement
condition.

To distinguish semantic predicates from their effective presentations,
consider also the maps
\[
 \operatorname{wp}_f,\operatorname{wlp}_f:
 \mathcal P(\N)\longrightarrow\mathcal P(\N)
\]
on the full powerset, defined by
\begin{align*}
 \operatorname{wp}_f(Q)
 &=\{n:f(n)\text{ is defined and belongs to }Q\},\\
 \operatorname{wlp}_f(Q)
 &=(\N\setminus\operatorname{dom}f)\cup\operatorname{wp}_f(Q).
\end{align*}
These satisfy
\begin{equation}\label{hp:halting-predicates}
 \operatorname{wp}_f(\N)=\operatorname{dom}f,
 \qquad
 \operatorname{wlp}_f(\varnothing)=\N\setminus\operatorname{dom}f.
\end{equation}
For recursive $Q$, membership in $\operatorname{wp}_f(Q)$ is c.e.,
whereas membership in $\operatorname{wlp}_f(Q)$ is co-c.e.:
its complement is $\operatorname{wp}_f(\N\setminus Q)$. Neither
predicate is generally recursive. If $Q$ is given by an arithmetic
formula, a formula for either transformed predicate can nevertheless be
constructed effectively from a program index, by coding finite
computations. Construction of such a formula does not decide its truth.

There is also no uniform procedure that, from an arbitrary index for a
partial recursive function with recursive domain, returns a c.e. index
for the complement of that domain. For a machine index $e$, let
$f_e(n)$ simulate machine $e$ on input $0$ and return $0$ if that
simulation halts. The domain of $f_e$ is either $\N$ or
$\varnothing$, so it is recursive in every case. A uniform index for
its complement would enumerate $0$ exactly when that simulation does
not halt. This would make the complement of the halting set c.e.
Pointwise existence in Theorem~\ref{hp:numerical-wlp} therefore does not
give uniform synthesis from program indices.

\section{Stable unions of domains}\label{ct:union-section}

Complementary domains and their inverse images enter the separation
arguments. This section assumes binary restriction products and the universal
range property of Di Paola--Heller \cite[Definition~5.1]{DPH}.
Explicitly, for
$f:X\to Y$, its range receives $f$ and satisfies
\begin{equation}\label{ct:universal-range}
 \beta f=\beta'f\quad\Longrightarrow\quad
 \beta\ran(f)=\beta'\ran(f)
 \qquad(\beta,\beta':Y\to Z).
\end{equation}
This is an additional hypothesis on the range structure.

Write $X\boxtimes Y$ for a binary restriction product. Its projections are
total, and the pairing of $f:Z\to X$ and $g:Z\to Y$ is characterized by
\[
 \pi_1\langle f,g\rangle=f\bar g,\qquad
 \pi_2\langle f,g\rangle=g\bar f,\qquad
 \overline{\langle f,g\rangle}=\bar f\bar g.
\]
The induced product of arrows is a bifunctor, with the usual canonical
associativity and symmetry isomorphisms. Finite powers below have fixed
parenthesizations. Total arrows form a category $\Tot(\R)$ with binary
products and diagonals $\Delta_X=\langle1_X,1_X\rangle$.

\begin{definition}[Stable union]\label{ct:stable-union}
Let $(e_i)_{i\in I}$ be a nonempty finite or countable family of domains on
$X$. A domain $e$ is its \emph{union} if $e_i\le e$ for every $i$ and
\[
 (\forall i,\ \beta e_i=\beta'e_i)
 \quad\Longrightarrow\quad \beta e=\beta'e
 \qquad(\beta,\beta':X\to Y).
\]
It is \emph{stable} if $1_W\boxtimes e$ is a union of
$(1_W\boxtimes e_i)_i$ for every object $W$.
\end{definition}

This is Heller's definition \cite[\S6, p.~1262]{Heller}. Taking $\beta$ to
be an upper-bound domain and $\beta'=1_X$ shows that a union is a least
upper bound, hence unique. Write it as $\bigvee_i e_i$ and put
$\bigvee\varnothing=0_X$.

\begin{lemma}\label{ct:union-calculus}
Under \eqref{ct:universal-range}, multiplication by a fixed domain and
direct image preserve unions. Inverse image along a total arrow preserves
stable unions. \end{lemma}
\begin{proof}
Agreement after all $de_i$ implies agreement after $de$ by applying the
union property to $\beta d,\beta'd$. For direct image, agreement after
all $\ran(fe_i)$ first implies agreement after all $fe_i$, then after
$fe$, and finally after $\ran(fe)$ by \eqref{ct:universal-range}.

For a total $h:X\to Y$, put $g=\langle1_X,h\rangle$ and $r=\ran(g)$.
For every domain $v$ on $Y$, the restriction and range identities give
\[
 (1_X\boxtimes v)r=\ran(g h^*v),\qquad
 (\pi_1)_!((1_X\boxtimes v)r)=h^*v.
\]
Indeed, $(1_X\boxtimes v)g=g h^*v$, and $\pi_1g=1_X$.
If $v=\bigvee_i v_i$ is stable, multiply the union of
$1_X\boxtimes v_i$ by $r$ and apply $(\pi_1)_!$.
The first two assertions give $h^*v=\bigvee_i h^*v_i$.
Applying the same argument to $1_W\boxtimes h$ proves stability of this
union.

\end{proof}

Assume finite stable unions exist. Each $\D(X)$ is then a bounded
distributive lattice. Let $\BDLat$ denote the category of such lattices
and maps preserving finite meets, finite joins, and both bounds. There is
a functor
\[
 E:\Tot(\R)^{\op}\longrightarrow\BDLat.
\]
It assigns $E(X)=\D(X)$ and $E(f^{\op})=f^*$. The restriction identities
give its functor laws and meet preservation; totality gives $f^*1=1$,
and Lemma~\ref{ct:union-calculus} gives join preservation.


\section{Categorical separation and proof codes}\label{ct:transfer-section}

Retain finite stable unions. A pair of domains $(e_+,e_-)$ on $X$ is
\emph{complementary} if $e_+e_-=0_X$ and $e_+\vee e_-=1_X$.
It \emph{separates} domains $a,b$ if $a\le e_+$ and $b\le e_-$.
Disjoint domains $a,b$ are \emph{inseparable} if no complementary pair
separates them. Below a domain $s$, the relative definitions replace
$1_X$ by $s$. This agrees with the complemented-domain separation
criterion of Di Paola--Heller on the ambient object
\cite[Definition 8.14]{DPH}. Relative inseparability implies ambient
inseparability: intersect an ambient complementary pair with $s$.

On an assertion domain $s$ on $U$, supply an arrow
\[
 \nu:U\to U,\qquad \bar\nu=s,\qquad s\nu=\nu.
\]
For $t\le s$, define $t^-=\nu^*t$. Call $t$ \emph{consistent} if
$tt^-=0_U$, and \emph{complete} if $t\vee t^-=s$.
In syntax, $\nu$ constructs the code of a negated assertion. Its
inverse-image action on domains is monotone.

\begin{theorem}[Incompleteness from inseparable domains]\label{ct:certificate-transfer}
Let $\R$ be a range restriction category with binary restriction
products, the universal range property, and finite stable unions as in
Section~\ref{ct:union-section}. Let $a,b\in\D(X)$ be inseparable.
Suppose there are an assertion domain $s\in\D(U)$, an arrow
$\nu:U\to U$ with $\bar\nu=s$ and $s\nu=\nu$, and a
proof-conclusion map $d_0:V\to U$ with $t_0=\ran(d_0)\le s$.
Suppose further that $\theta:X\to U$ is total, $s\theta=\theta$,
and that arrows
\[
 \pi_+:W_+\to X,\quad \pi_-:W_-\to X,\qquad
 u_+:W_+\to V,\quad u_-:W_-\to V
\]
satisfy
\begin{equation}\label{ct:certificate-equations}
 \ran(\pi_+)=a,\quad \ran(\pi_-)=b,\qquad
 d_0u_+=\theta\pi_+,\quad d_0u_-=\nu\theta\pi_-.
\end{equation}
Every consistent domain $t$ with $t_0\le t\le s$ is incomplete.
\end{theorem}
\begin{proof}
The first commuting equation and $t_0d_0=d_0$ give
$\ran(\theta\pi_+)\le t_0$. The range axiom and the local adjunction
for the total arrow $\theta$ give
\[
 \theta_!a=\ran(\theta\ran(\pi_+))
 =\ran(\theta\pi_+)\le t_0,
 \qquad a\le\theta^*t_0.
\]
Similarly $\nu\theta$ is total, since $\bar\nu=s$ and
$s\theta=\theta$, and the second commuting equation gives
$b\le\theta^*\nu^*t_0$. For an extension $t$, monotonicity yields
\[
 a\le\theta^*t,\qquad b\le\theta^*\nu^*t.
\]
If $t$ were consistent and complete, preservation of meets and finite
stable unions by total inverse image would make these two right-hand
domains complementary on $X$, because $\theta^*s=1_X$. They would
separate $a,b$, contrary to inseparability.
\end{proof}

In a deductive presentation, $W_+$ and $W_-$ can be objects of finite
computation evidence, and $u_+,u_-$ construct proofs of the two
conclusions. Their projections supply the inseparable domains by taking
ranges. The theorem applies to deductively closed extensions when their
theorem domains are among the domains $t$ in its statement.

\subsection{A representation condition implying inseparability}

There is a second route in which inseparability follows from a
representation property. Fix a nontrivial object $U$, meaning
$0_U\ne1_U$, a domain $s$ on $U$, disjoint domains $p_+,p_-\le s$,
and a total arrow
\[
 \sigma:U\boxtimes U\to U,\qquad s\sigma=\sigma.
\]
For every total $g:U\to U$, put
$f_g=\sigma(g\boxtimes1_U):U\boxtimes U\to U$.
Assume that for every complementary pair $(e_+,e_-)$ on $U$ there is
such a $g$ with
\begin{equation}\label{ct:representation}
 f_g^*p_+=\pi_2^*e_+,\qquad
 f_g^*p_-=\pi_2^*e_-.
\end{equation}

\begin{theorem}\label{ct:representation-inseparable}
Under these hypotheses, $p_+,p_-$ are inseparable relative to $s$.
\end{theorem}
\begin{proof}
Suppose a complementary pair $(d_+,d_-)$ relative to $s$ separates them.
Put $\delta=\sigma\Delta_U$. Total inverse image makes
$e_+=\delta^*d_-$ and $e_-=\delta^*d_+$ complementary on $U$.
Choose $g$ satisfying \eqref{ct:representation}, and put
$\gamma=\langle1_U,g\rangle$ and $h=\delta g$. Then
\[
 f_g\gamma=h,\qquad\pi_2\gamma=g.
\]
The first identity gives the commutative square
\[
\begin{tikzcd}[column sep=large]
 U \arrow[r,"\gamma"] \arrow[d,"g"'] &
 U\boxtimes U \arrow[d,"f_g"]\\
 U \arrow[r,"\delta"'] & U.
\end{tikzcd}
\]
Pullback along $\gamma$ gives
$h^*p_+=h^*d_-$ and $h^*p_-=h^*d_+$.
The separator inequalities imply both
$h^*d_-\le h^*d_+$ and $h^*d_+\le h^*d_-$.
These two domains are complementary, because $h$ is total and
$sh=h$. Equality therefore forces $0_U=1_U$, a contradiction.
\end{proof}

For $p_+=b$ and $p_-=\nu^*b$, every consistent domain extension
$b\le t\le s$ is consequently incomplete: a complete extension would
itself be a relative complementary separator. The representation
condition \eqref{ct:representation} supplies no uniform assignment
of $g$ to a code of $(e_+,e_-)$. Its conclusion is inseparability, as
defined above; effective inseparability requires further uniform data.

\section{The numerical category and effective unions}
\label{num:ambient}

Let $\R_{\mathrm{rec}}$ have objects $\N^k$, for $k\ge1$, and partial
recursive maps as arrows, with equality of partial functions. Composition
is composition of partial maps. Restriction sends a map to the partial
identity on its domain. The restriction product is the Cartesian product
of underlying sets, acting on pairs of partial maps where both components
are defined. Recursive pairing gives isomorphisms $\N^k\cong\N$.
A universal partial recursive function supplies a Turing morphism, making
this the numerical recursion category of \cite{DPH,Heller}.

A set is c.e. precisely when it is the domain of a partial recursive
function. Its partial identity is therefore an arrow. Conversely, the
domain of every restriction idempotent is c.e. For a c.e. subset
$A\subseteq\N^k$, write $1_A$ for this partial identity. The range of
$f$ is $1_{\operatorname{im}f}$, since the image of a partial recursive
map is c.e. It has the universal equality property
\eqref{ct:universal-range}: agreement after $f$ is agreement, including
definedness, at every point of its image.

\begin{proposition}\label{num:uniform-unions}
Suppose $(A_j)_{j\in\N}$ is a uniformly c.e. family of subsets of
$\N^k$, meaning that $\{(j,x):x\in A_j\}$ is c.e. Then
$1_{\bigcup_j A_j}$ is the stable union of the domains $1_{A_j}$ in
$\R_{\mathrm{rec}}$.
\end{proposition}
\begin{proof}
Dovetailing the enumerations gives a c.e. union. If two partial maps agree
after every $1_{A_j}$, they have the same definedness and value at each
point of the union. This proves Heller's universal agreement condition.
For every $m\ge1$, the same argument applies to
$\N^m\times\bigcup_j A_j=\bigcup_j(\N^m\times A_j)$, proving stability.
\end{proof}

The uniformity condition is substantive. Let $B\subseteq\N$ be a set
that is not c.e., and put $B_j=B\cap\{0,\ldots,j\}$. Each $B_j$ is
finite, but the increasing family has no supremum among c.e. sets.
Indeed, a supremum would contain $B$, while every
$\N\setminus\{x\}$ with $x\notin B$ is an upper bound; the supremum
would therefore be $B$.


\section{A Turing morphism and inseparable domains}
\label{tur:section}

A Turing morphism supplies the computational obstruction used by the
separation theorem. We work in a restriction category with zero arrows,
binary restriction products, and finite joins of domains preserved by
inverse image along total arrows. The restriction-product conventions are
those of Section~\ref{ct:union-section}. The domain hypotheses hold, in
particular, under the stable-union hypotheses of
Lemma~\ref{ct:union-calculus}.

Fix an object $U$ with $0_U\ne1_U$ and an arrow
\[
 \tau:U\boxtimes U\longrightarrow U.
\]
The instance of Turing universality needed below is the following:
for every arrow $F:U\boxtimes U\to U$, there is a total arrow
$g:U\to U$ such that
\begin{equation}\label{tur:universality}
 \tau(g\boxtimes1_U)=F.
\end{equation}
The arrow $\tau$ may be partial. The total arrow $g$ assigns a program
parameter; evaluation of that parameter remains subject to the domain of
$\tau$.

Choose disjoint domains $q_0,q_1$ on $U$. We impose a separate condition
on binary case distinction: whenever $c,c'$ are complementary domains on
$U$, there is a total arrow $h:U\to U$ satisfying
\begin{equation}\label{tur:switch}
 h^*q_0=c',\qquad h^*q_1=c.
\end{equation}
This condition constructs an arrow from a domain decomposition. For
example, suppose compatible arrows have joins preserved by composition,
restriction preserves these joins, and there are total arrows
$k_i:U\to U$ with $q_i k_i=k_i$, for $i=0,1$. Here arrows $f,g$ with
common source and target are \emph{compatible} when
$f\bar g=g\bar f$. The disjoint restrictions $c,c'$ make $k_1c$ and
$k_0c'$ compatible, and
\[
 h=k_1c\vee k_0c'
\]
is total. Disjointness of $q_0,q_1$ gives
$q_0h=k_0c'$ and $q_1h=k_1c$, which prove
\eqref{tur:switch}. Thus binary case distinction follows from suitable
joins of arrows, in addition to the joins of domains.

\begin{theorem}\label{tur:inseparable}
Let $\R$ be a restriction category with zeros, binary restriction
products, and finite joins of domains preserved by total inverse image.
Suppose $U$ is nontrivial, $\tau:U\boxtimes U\to U$ satisfies
\eqref{tur:universality}, and disjoint domains $q_0,q_1$ satisfy the
binary case condition \eqref{tur:switch}. Put
\[
 d=\tau\Delta_U:U\longrightarrow U,\qquad
 a=d^*q_0,\qquad b=d^*q_1.
\]
Then $a,b$ are disjoint and inseparable.
\end{theorem}
\begin{proof}
Inverse image preserves meets and zero, so
$ab=d^*(q_0q_1)=0_U$. Suppose complementary domains $c,c'$ separate
$a,b$, with $a\le c$ and $b\le c'$. Choose $h$ as in
\eqref{tur:switch}. Apply \eqref{tur:universality} to
$F=h\pi_2$ to obtain a total $g:U\to U$ with
\[
 \tau(g\boxtimes1_U)=h\pi_2.
\]
Since $g$ is total, the product identities give
\[
 (g\boxtimes1_U)\langle1_U,g\rangle=\Delta_U g,
 \qquad \pi_2\langle1_U,g\rangle=g.
\]
Precomposition therefore yields
\begin{equation}\label{tur:diagonal-equality}
 dg=hg.
\end{equation}
The corresponding commutative square is
\[
\begin{tikzcd}[column sep=large]
 U \arrow[r,"{\langle1_U,g\rangle}"] \arrow[d,"g"'] &
 U\boxtimes U \arrow[d,"\tau(g\boxtimes1_U)=h\pi_2"]\\
 U \arrow[r,"d"'] & U.
\end{tikzcd}
\]
Pulling back $q_0,q_1$ along \eqref{tur:diagonal-equality} gives
\[
 g^*a=g^*h^*q_0=g^*c',\qquad
 g^*b=g^*h^*q_1=g^*c.
\]
The separator inequalities consequently imply
$g^*c'\le g^*c$ and $g^*c\le g^*c'$. Totality of $g$ makes
$g^*c,g^*c'$ complementary. Their equality forces $0_U=1_U$,
a contradiction.
\end{proof}

The universality of $\tau$ is used to construct $g$ and hence the
identity \eqref{tur:diagonal-equality}. The categorical diagonal enters
as a total structural arrow. The proof requires neither a global element
of $U$ nor a choice of a global program code.

\begin{proposition}[Numerical realization]\label{tur:numerical}
In $\R_{\mathrm{rec}}$, let $\tau:\N^2\to\N$ be a universal partial
recursive evaluator for an acceptable numbering, and let
$q_i=1_{\{i\}}$, for $i=0,1$. Then
\[
 A=\{n:\tau(n,n)=0\},\qquad
 B=\{n:\tau(n,n)=1\}
\]
are disjoint c.e. sets whose domains $1_A,1_B$ are inseparable.
\end{proposition}
\begin{proof}
Parameterization of partial recursive functions gives
\eqref{tur:universality}. For complementary c.e. sets $C,C'$, dovetail
their membership semidecisions. Exactly one succeeds on every input.
Output $1$ when the test for $C$ succeeds and $0$ when the test for
$C'$ succeeds. This defines a total recursive $h$ satisfying
\eqref{tur:switch}. Finite unions of c.e. sets are c.e., and recursive
inverse image preserves these unions. Theorem~\ref{tur:inseparable}
therefore applies. A complementary pair of c.e. domains is the pair of
a recursive set and its complement, so inseparability here is ordinary
recursive inseparability.
\end{proof}

\section{Typed inference rules and polynomial functors}\label{poly:typed}

An inference rule has a conclusion and a family of premises. Polynomial
functors record this dependence before a deductive closure is taken.
Let $S$ be a set of expressions. A typed rule presentation consists of a
set $A$ of rule instances, a finite set $B_a$ of premise positions for
each $a\in A$, and maps
\[
 t:A\longrightarrow S,\qquad
 s:\coprod_{a\in A}B_a\longrightarrow S.
\]
The expression $t(a)$ is the conclusion of $a$, and $s(a,b)$ is its
premise at position $b$. A rule with $B_a=\varnothing$ is an axiom
instance. Regard $S$ as a discrete category, so that $\Set^S$ is the
category of $S$-indexed families of sets and componentwise maps.

Define the functor
\begin{equation}\label{poly:family-type}
 P:\Set^S\longrightarrow\Set^S
\end{equation}
on objects by
\begin{equation}\label{poly:family-formula}
 (PX)_\sigma=
 \coprod_{\substack{a\in A\\t(a)=\sigma}}
 \prod_{b\in B_a}X_{s(a,b)}.
\end{equation}
For a family of maps $v_\sigma:X_\sigma\to Y_\sigma$, put
\[
 (Pv)_\sigma(a,(x_b)_{b\in B_a})
   =(a,(v_{s(a,b)}(x_b))_{b\in B_a}).
\]
The identity and composition laws hold componentwise. Thus $PX$ is the
family of one-step deductions whose premise proofs are drawn from $X$.

This presentation is also a polynomial bicomodule. Write $\mathbf{Poly}$
for the category of polynomial endofunctors of $\Set$ and natural
transformations. Functor composition gives its monoidal product
$p\mathbin{\triangleleft}q=p\circ q$, with unit
$y=\id_{\Set}$; the associativity and unit constraints are those of
functor composition \cite[\S6.1.1]{NS2024}. Write $y^B(X)=X^B$ and
$Sy=\coprod_{s\in S}y$. The discrete category $S$ corresponds to the
polynomial comonoid $Sy$, whose counit forgets the
$S$-label and whose comultiplication repeats it. The rule polynomial is
\[
 m=\coprod_{a\in A}y^{B_a}.
\]
Its left coaction $\lambda:m\to(Sy)\triangleleft m$ sends the position
$a$ to $(t(a),a)$ and is the identity on premise positions. Its right
coaction $\rho:m\to m\triangleleft(Sy)$ sends $a$ to
$(a,b\mapsto s(a,b))$ and is again the identity on premise positions.
Repeating or forgetting these labels proves the two coaction laws.
The two orders of labelling premises and conclusions give the same
position and the same direction map, so
\[
 ((Sy)\triangleleft\rho)\lambda
   =(\lambda\triangleleft(Sy))\rho.
\]
Thus $m$ is an $(S,S)$-bicomodule in the sense of
\cite[Definition~8.96]{NS2024}. Its associated functor is
\eqref{poly:family-type}; compare
\cite[Proposition~8.104]{NS2024}.

The premise direction in a polynomial morphism is contravariant.
For two presentations over $S$, a morphism of their bicomodules is a
map $u:A\to A'$ together with maps
\[
 \beta_a:B'_{u(a)}\longrightarrow B_a
\]
such that
\[
 t'(u(a))=t(a),\qquad
 s'(u(a),b')=s(a,\beta_a(b')).
\]
The corresponding natural transformation $P\Rightarrow P'$ sends
$(a,(x_b))$ to $(u(a),(x_{\beta_a(b')}))$. Hence a target premise may
reuse a source premise, and a source premise may be omitted. A
cartesian polynomial morphism has bijective direction maps
\cite[Definition~5.51]{NS2024}.

\section{Finite proofs in a computational category}
\label{poly:computational}

Splitting restriction idempotents makes the domains of computations
available as objects. For a restriction category $\R$, define
$\operatorname{Split}_r(\R)$ as follows. Its objects are pairs $(X,e)$,
where $e\in\D(X)$. An arrow
$f:(X,e)\to(Y,e')$ is an arrow $f:X\to Y$ of $\R$ satisfying
$f=e'fe$; its restriction is $\bar f$, and the identity of $(X,e)$ is
$e$. Composition is inherited from $\R$. Such an arrow is total
precisely when $\bar f=e$.

For the numerical category of Section~\ref{num:ambient}, put
\[
 \mathscr C=
 \Tot\bigl(\operatorname{Split}_r(\R_{\mathrm{rec}})\bigr).
\]
Concretely, an object is a c.e. subset of a finite power of $\N$, and
an arrow is a computable function defined on its entire source and
taking values in its target. Two arrows are equal when their values
agree on the source. Finite products, pullbacks, and finite coproducts
exist in $\mathscr C$: tuples give products, equality of computed
values gives pullbacks, and recursive tags give coproducts. Each
resulting carrier is c.e. A singleton is terminal and the empty set is
initial.

Domains on objects of $\mathscr C$ will mean restriction idempotents
in $\operatorname{Split}_r(\R_{\mathrm{rec}})$. On a c.e. carrier
$X$, they are the partial identities on c.e. subsets of $X$.
In the splitting, ranges and binary restriction products are computed
on these c.e. carriers. Finite unions satisfy Heller's agreement and
stability conditions by the pointwise argument of
Proposition~\ref{num:uniform-unions}, now allowing arbitrary c.e.
carriers. Ranges have the same universal equality property.

Fix a recursive set $S$ of expression codes. A \emph{recursive rule
presentation} $\Sigma$ consists of recursive sets $A_0,A_1,A_2$ and
computable maps
\[
 s_i:A_i\longrightarrow S^i,\qquad
 t_i:A_i\longrightarrow S\qquad(0\le i\le2).
\]
The integer $i$ is the number of ordered premises; $S^0$ is the
terminal object. All these sets and maps are objects and arrows of
$\mathscr C$. A c.e. list of rules with at most two premises can be
presented in this form by
including a finite enumeration witness in each rule instance, provided
the premise and conclusion maps are computable from that witnessed
instance.

Define the functor
\begin{equation}\label{poly:slice-type}
 P_\Sigma:\mathscr C/S\longrightarrow\mathscr C/S.
\end{equation}
For $q:X\to S$, let its value be
\begin{equation}\label{poly:slice-formula}
 P_\Sigma(q)=
 \left(
 \coprod_{i=0}^{2}A_i\times_{S^i}X^i
 \xrightarrow{[t_i\pi_{A_i}]}S
 \right),
\end{equation}
where the pullback uses $s_i$ and $q^i$. For a slice arrow
$v:(X,q)\to(Y,r)$, the map $P_\Sigma(v)$ fixes the rule instance and
applies $v$ to each premise entry. It is computable and total on the
displayed source. The identity and composition laws follow by applying
the corresponding laws for $v$ to each premise.

\begin{theorem}[Finite proof algebra]\label{poly:initial}
For every recursive rule presentation $\Sigma$ over a recursive
expression object $S$, the functor $P_\Sigma$ of
\eqref{poly:slice-type} has an initial algebra
\[
 \iota_\Sigma:P_\Sigma(c_\Sigma)\longrightarrow c_\Sigma
 \quad\text{in }\mathscr C/S.
\]
Here $D_\Sigma$ is the recursive set of valid finite ordered
rule-labelled proof trees, $c_\Sigma:D_\Sigma\to S$ gives the root
conclusion, and $\iota_\Sigma$ grafts a rule instance onto its premise
proofs.
\end{theorem}
\begin{proof}
Fix a recursive coding of finite ordered labelled trees. To check a
code, inspect its finitely many nodes. A node labelled by $a\in A_i$
must have exactly $i$ children. The tuple of their conclusions must be
$s_i(a)$, and the conclusion of the node is $t_i(a)$. Membership in
$A_i$ is decidable; after it holds, all required values are computable.
This is a terminating decision procedure for $D_\Sigma$. The root
conclusion and grafting operations are computable on their domains.

Let $\alpha:P_\Sigma(q)\to q$ be an algebra, where $q:X\to S$ is an
arbitrary object of $\mathscr C/S$. Define $h:D_\Sigma\to X$ by
evaluating each finite tree from its leaves to its root using
$\alpha$. At a leaf, its zero-premise rule is in the domain of
$\alpha$. Inductively, the evaluated children belong to $X$ and have
the prescribed conclusions, so they and the rule instance form an
element of the pullback in \eqref{poly:slice-formula}. The next
application of $\alpha$ is therefore defined. This gives a
computable function total on $D_\Sigma$, even when $X$ is only c.e.
The same induction proves
\[
 qh=c_\Sigma,\qquad
 h\iota_\Sigma=\alpha P_\Sigma(h).
\]
Thus $h$ is an algebra morphism in the slice. These equations determine
its value at every leaf and then at every larger finite tree, proving
uniqueness.
\end{proof}

The initiality assertion takes place in a category of computable maps:
the fold into every algebra must itself be computable. It is distinct
from the cofree polynomial comonoid of
\cite[Theorem~8.45]{NS2024}, whose objects include infinite behaviour
trees. For example, the cofree comonoid on $2y$ already has uncountably
many objects \cite[Example~8.10]{NS2024}; it cannot serve as the
recursive object of finite proofs constructed here.

\subsection{Translations of presentations}\label{poly:translations}

Let $\mathbf{Pres}_S$ be the category whose objects are recursive rule
presentations over $S$. An arrow $u:\Sigma\to\Sigma'$ consists of
computable maps $u_i:A_i\to A_i'$ satisfying
\[
 s_i'u_i=s_i,\qquad t_i'u_i=t_i.
\]
Identities and composition are componentwise. Such a translation
preserves each premise position and induces a natural transformation
$\eta_u:P_\Sigma\Rightarrow P_{\Sigma'}$ by replacing the rule label
$a$ with $u_i(a)$. For a slice map $v$, both composites in its
naturality square apply $u_i$ to the label and $v$ to each premise.

Define the functor
\[
 \mathsf{Proof}_S:\mathbf{Pres}_S\longrightarrow\mathscr C/S
\]
on objects by $\mathsf{Proof}_S(\Sigma)=c_\Sigma$. On an arrow $u$,
initiality gives the unique $P_\Sigma$-algebra morphism
\[
 D(u):D_\Sigma\longrightarrow D_{\Sigma'}
\]
into the algebra
$\iota_{\Sigma'}\eta_{u,c_{\Sigma'}}$. Equivalently, $D(u)$ replaces
every rule label in a finite proof by its image under $u$. It satisfies
\[
 \begin{tikzcd}
 D_\Sigma\arrow[r,"D(u)"]\arrow[d,"c_\Sigma"']
   &D_{\Sigma'}\arrow[d,"c_{\Sigma'}"]\\
 S\arrow[r,equal]&S.
 \end{tikzcd}
\]
Both $D(1_\Sigma)$ and $1_{D_\Sigma}$ are folds into the same algebra,
so they agree. Likewise, uniqueness gives $D(vu)=D(v)D(u)$.
This proves functoriality. Literal additions of rule instances give
arrows of $\mathbf{Pres}_S$.

\subsection{Ranges of conclusion maps}\label{poly:ranges}

The passage from proof objects to theorem objects uses only ranges.
The following construction therefore applies beyond the numerical
category.

\begin{proposition}\label{poly:range-functor}
Let $\R$ be a range restriction category and put
$\mathscr E=\Tot(\operatorname{Split}_r(\R))$. There is a functor
\[
 \mathsf{Ran}:\mathscr E^{[1]}\longrightarrow\mathscr E
\]
and a natural monomorphism
$j:\mathsf{Ran}\Rightarrow\mathsf{Cod}$, where $[1]$ is the category
with one nonidentity arrow and $\mathsf{Cod}$ takes an arrow to its
codomain. For $c:(Y,d)\to(X,e)$, its value is
$\mathsf{Ran}(c)=(X,\ran(c))$ and $j_c$ is the range inclusion.
\end{proposition}
\begin{proof}
Write $t=\ran(c)$. Since $ec=c$, the least-domain property of ranges
gives $t\le e$. Thus $(X,t)$ is an object and the underlying arrow
$t:(X,t)\to(X,e)$ is total. It is monic: for every arrow $v$ into
$(X,t)$, the equation $tv=v$ holds.

An arrow of $\mathscr E^{[1]}$ is a commuting square
\[
 \begin{tikzcd}
 (Y,d)\arrow[r,"h"]\arrow[d,"c"']
   &(Y',d')\arrow[d,"c'"]\\
 (X,e)\arrow[r,"f"']&(X',e').
 \end{tikzcd}
\]
Put $t'=\ran(c')$. The range axioms give
\[
 \ran(ft)=\ran(fc)=\ran(c'h)\le t',
\]
because $t'c'h=c'h$. Consequently $t'ft=ft$. Moreover,
$\overline{ft}=\bar f t=et=t$, since $f$ is total from $(X,e)$.
It follows that $ft:(X,t)\to(X',t')$ is a total arrow; define
$\mathsf{Ran}(h,f)=ft$. Identities map to the corresponding range
idempotents. If the next square has bottom arrow $g$ and middle
range $t'$, then
$(gt')(ft)=gft$, proving the composition law. Finally,
$j_{c'}\mathsf{Ran}(h,f)=ft=fj_c$, which proves naturality.
\end{proof}

Apply this proposition to the computational proof functor. The theorem
domain of $\Sigma$ is
\[
 t_\Sigma=\ran(c_\Sigma).
\]
It is c.e., since it is the image of a computable map on a recursive
set. For $u:\Sigma\to\Sigma'$, the equality
$c_{\Sigma'}D(u)=c_\Sigma$ implies
$t_\Sigma\le t_{\Sigma'}$. Thus the range inclusions define a functor
\[
 \mathsf{Th}_S:\mathbf{Pres}_S\longrightarrow\mathscr C/S,
 \qquad
 \Sigma\longmapsto\bigl(j_{c_\Sigma}:
       \mathsf{Ran}(c_\Sigma)\longrightarrow S\bigr).
\]
On arrows the functor is the inclusion specified by
$t_\Sigma\le t_{\Sigma'}$. Forgetting
the maps to $S$ identifies this functor with the composite of
$\mathsf{Proof}_S$, the inclusion of the slice into the arrow category,
and $\mathsf{Ran}$.

\section{Profunctors and polynomial bicomodules}\label{poly:hoare}

The polynomial description of rules and the Hoare profunctors admit a
comparison through functor categories. Let $A,B$ be small categories and
let $H:A^{\op}\times B\to\Set$ be a profunctor. Write $\Set^A$ and
$\Set^B$ for the categories of covariant set-valued functors. Define
\[
 R_H:\Set^B\longrightarrow\Set^A
\]
on objects by
\[
 (R_HX)(a)=\operatorname{Nat}(H(a,-),X).
\]
For $f:a\to a'$ in $A$, its action is precomposition with
$H(f^{\op},1):H(a',-)\Rightarrow H(a,-)$. For a natural transformation
$\xi:X\Rightarrow X'$, the component $(R_H\xi)_a$ is
postcomposition with $\xi$. Associativity and the identity laws for
composition of natural transformations give both functor laws.

\begin{proposition}\label{poly:profunctor-adjunction}
For every profunctor $H:A^{\op}\times B\to\Set$ between small
categories, $R_H$ has a left adjoint
\[
 L_H:\Set^A\longrightarrow\Set^B,\qquad
 (L_HY)(b)=\int^{a\in A}Y(a)\times H(a,b).
\]
Consequently $R_H$ preserves all small limits. For a profunctor
$K:B^{\op}\times C\to\Set$ there is a natural isomorphism
\[
 R_{K\odot H}\cong R_H R_K.
\]
\end{proposition}
\begin{proof}
For $g:b\to b'$, the map $(L_HY)(g)$ sends $[y,h]$ to
$[y,H(1,g)h]$. For $\eta:Y\Rightarrow Y'$, the map
$(L_H\eta)_b$ sends $[y,h]$ to $[\eta_a(y),h]$. The coend
relations and naturality show that these maps are well-defined and
satisfy the functor laws. A natural transformation $L_HY\Rightarrow X$
is equivalently a family of maps
$Y(a)\times H(a,b)\to X(b)$, natural in $b$ and compatible with the
coend identifications in $a$. Currying gives exactly a natural
transformation $Y\Rightarrow R_HX$. This bijection is natural in
$Y,X$, so $L_H\dashv R_H$.

For the composition statement, the coend universal property and currying
give, naturally in $a$ and $X$,
\[
 \operatorname{Nat}\bigl((K\odot H)(a,-),X\bigr)
 \cong
 \operatorname{Nat}\bigl(H(a,-),R_KX\bigr).
\]
These are the components of the asserted isomorphism.
\end{proof}

Niu and Spivak identify $(A,B)$-bicomodules with parametric right adjoints
$\Set^B\to\Set^A$ \cite[Proposition~8.104, p.~336]{NS2024}.
Since $R_H$ is a right adjoint, it preserves connected limits and so
determines such a bicomodule. The functor $R_H$ has one operation at each
object $a$, whose input functor is $H(a,-)$.

For the Hoare profunctors of Section~\ref{hp:profunctors}, Yoneda gives
\[
 (R_{K_f}X)(p)\cong X(f_!p),\qquad
 R_{H_f}\cong\operatorname{Ran}_{f^*}.
\]
Here the right Kan extension has type
$\Set^{\D(Y)}\to\Set^{\D(X)}$ and is right adjoint to
precomposition along $f^*:\D(Y)\to\D(X)$. Indeed,
$L_{H_f}Y\cong Y\circ f^*$ by co-Yoneda, while the first displayed
isomorphism follows from
$K_f(p,-)=\D(Y)(f_!p,-)$.
These adjoints act on set-valued functors. A weakest liberal precondition
in Section~\ref{hp:representability} instead represents a presheaf by
a domain in $\D(X)$; its existence is the separate condition tested by
Theorem~\ref{hp:numerical-wlp}.

\section{A numerical deductive construction}
\label{num:section}

Finite computation traces determine nullary inference rules. The polynomial
proof construction then supplies conclusion maps whose ranges contain the
two computational domains in opposite logical polarities.

\subsection{Axioms, proofs, and refutations}

Fix the universal partial recursive evaluator
$\tau:\N^2\rightharpoonup\N$, and put
\[
 d(n)=\tau(n,n),\qquad
 A=\{n:d(n)=0\},\qquad B=\{n:d(n)=1\}.
\]
The numerical instance of Theorem~\ref{tur:inseparable} gives
inseparability of the partial identities $a=1_A$ and $b=1_B$.
Equivalently, no recursive set contains $A$ and is disjoint from $B$:
a complementary pair of c.e. subsets of $\N$ is decidable by dovetailing
its two enumerations.

Fix a machine computing $d$. For $i=0,1$, let
\[
 W_i=\{(n,w):w\text{ is a finite halting computation of }d(n)
                    \text{ with output }i\}.
\]
Each $W_i$ is recursive: a proposed trace can be checked against the
machine's transition rules in finitely many steps. The projections
$\pi_i:W_i\to\N$, $(n,w)\mapsto n$, are arrows of $\mathscr C$, and
\[
 \ran(\pi_0)=a,\qquad \ran(\pi_1)=b.
\]

Let $S$ be the recursive set of codes for formulas generated by
\[
 \alpha::=\bot\mid\top\mid P_n\mid\neg\alpha
     \mid(\alpha\land\alpha)\mid(\alpha\lor\alpha),
 \qquad n\in\N.
\]
Write $\alpha\Rightarrow\beta$ for $\neg\alpha\lor\beta$.
The constructors and well-formedness test are recursive. In particular,
there are total computable maps
\[
 \theta:\N\longrightarrow S,\quad n\longmapsto\ulcorner P_n\urcorner,
 \qquad
 \nu:S\longrightarrow S,\quad
 \ulcorner\alpha\urcorner\longmapsto\ulcorner\neg\alpha\urcorner.
\]
We identify formulas with their codes when describing rule instances.

Define the presentation $\Sigma_0$ over $S$ as follows. Its nullary
instances form the tagged disjoint union
\[
 A_0=W_0\amalg W_1\amalg\mathsf{Taut},
\]
where $\mathsf{Taut}\subseteq S$ consists of classical propositional
tautologies. Their conclusions are, respectively, $P_n$, $\neg P_n$,
and the tautology itself. Truth-table evaluation decides membership in
$\mathsf{Taut}$. There are no unary instances. A binary instance is a
pair $(\alpha,\beta)\in S^2$ with ordered premises
$(\alpha,\alpha\Rightarrow\beta)$ and conclusion $\beta$.
Thus all instance sets and all premise and conclusion maps meet the
hypotheses of Theorem~\ref{poly:initial}.

Let
\[
 c^+:D\longrightarrow S
\]
be the underlying object in $\mathscr C/S$ of the initial algebra of
finite proofs, with $c^+$ the root conclusion map. Its range
$t^+=\ran(c^+)$ is the domain of theorems.
The underlying theorem set is
\begin{equation}\label{num:axioms}
 T_0=\operatorname{Cn}(\Gamma_0),\qquad
 \Gamma_0=\{P_n:n\in A\}\cup\{\neg P_n:n\in B\},
\end{equation}
where $\operatorname{Cn}$ means closure under tautologies and modus
ponens. The trace at each nonlogical axiom node is part of its proof code.

A refutation of $\alpha$ is a proof whose conclusion is $\neg\alpha$.
Define its object and conclusion map by the pullback in $\mathscr C$
\[
\begin{tikzcd}
 D^- \arrow[r,"r"] \arrow[d,"c^-"'] & D \arrow[d,"c^+"]\\
 S \arrow[r,"\nu"'] & S.
\end{tikzcd}
\]
Concretely,
\[
 D^-=\{(p,\alpha)\in D\times S:c^+(p)=\nu(\alpha)\},\qquad
 r(p,\alpha)=p,\qquad c^-(p,\alpha)=\alpha.
\]
The equality test is recursive, so $D^-$ is recursive and these maps are
computable. Its range satisfies
\begin{equation}\label{num:refutation-range}
 t^-:=\ran(c^-)=\nu^*t^+.
\end{equation}
Indeed, $\alpha$ lies in that range exactly when a proof of
$\nu(\alpha)$ exists.

Let $u_0:W_0\to D$ assign the one-node proof with conclusion $P_n$.
Let $u_1:W_1\to D^-$ assign the pair consisting of the one-node proof
of $\neg P_n$ and the formula $P_n$. These are total computable maps.
Their definitions give the commutative squares
\begin{equation}\label{num:certificate-maps}
\begin{tikzcd}[column sep=large]
 W_0 \arrow[r,"u_0"] \arrow[d,"\pi_0"'] & D \arrow[d,"c^+"]\\
 \N \arrow[r,"\theta"'] & S
\end{tikzcd}
\qquad
\begin{tikzcd}[column sep=large]
 W_1 \arrow[r,"u_1"] \arrow[d,"\pi_1"'] & D^- \arrow[d,"c^-"]\\
 \N \arrow[r,"\theta"'] & S.
\end{tikzcd}
\end{equation}
Applying the range construction of Proposition~\ref{poly:range-functor}
gives
\begin{equation}\label{num:range-inclusions}
 a\le\theta^*t^+,\qquad b\le\theta^*t^-.
\end{equation}
For example, $t^+\theta\pi_0=\theta\pi_0$ implies that the range of
$\pi_0$ is contained in $\theta^*t^+$; the other square gives the same
argument for $b$.
The negative square and the refutation pullback give
$c^+r u_1=\nu\theta\pi_1$. Hence these maps also realize the two
proof equations in Theorem~\ref{ct:certificate-transfer}, in the
numerical restriction-idempotent splitting.

\begin{lemma}\label{num:consistent}
The theorem domain $t^+$ is consistent: $t^+t^-=0_S$.
\end{lemma}
\begin{proof}
Give $P_n$ the truth value of $n\in A$ and extend this assignment by
the Boolean operations. Every axiom in $\Gamma_0$ is true, since
$A\cap B=\varnothing$. Tautologies are true, and modus ponens preserves
truth. Induction on the proof tree shows that every theorem is true.
A formula and its negation cannot both be true, proving the assertion.
\end{proof}

\begin{proposition}\label{num:proof-inseparable}
The domains $t^+,t^-$ on $S$ are inseparable.
\end{proposition}
\begin{proof}
They are disjoint by Lemma~\ref{num:consistent}. A complementary
separator for them would pull back along the total map $\theta$ to a
complementary separator for $a,b$, by \eqref{num:range-inclusions}.
\end{proof}

\subsection{Theories, deductive categories, and extensions}

A \emph{theory} in this language is a set $T\subseteq S$ containing
all tautologies and closed under modus ponens. An \emph{extension} of
$T$ is a theory $T'$ with $T\subseteq T'$. A theory is
\emph{consistent} if $\bot\notin T$, and \emph{complete} if every
formula $\alpha$ satisfies $\alpha\in T$ or $\neg\alpha\in T$.
The tautologies of classical logic imply that consistency is equivalent
to the absence of a formula whose two polarities both belong to $T$.

For each theory $T$, define a category $\mathcal B_T$ with formula
objects and hom sets
\[
 \mathcal B_T(\alpha,\beta)=
 \begin{cases}
 \{*\},&\alpha\Rightarrow\beta\in T,\\
 \varnothing,&\text{otherwise}.
 \end{cases}
\]
Reflexivity and transitivity tautologies give identities and composition;
the category laws hold because each hom set has at most one member.
Conjunction is a product, disjunction a coproduct, $\top$ terminal,
and $\bot$ initial, by their propositional universal properties.
Define the contravariant functor
\[
 N_T:\mathcal B_T^{\op}\longrightarrow\mathcal B_T
\]
by $N_T(\alpha)=\neg\alpha$ and by sending
$a^{\op}:\beta\to\alpha$ to the arrow
$\neg\beta\to\neg\alpha$ supplied by contraposition.
The functor laws follow from uniqueness of arrows. Double negation gives
a natural isomorphism
$1_{\mathcal B_T}\cong N_TN_T^{\op}$.
Thus negation induces an order-reversing involution on the Lindenbaum
Boolean algebra obtained by identifying mutually implying formulas.

For an inclusion $T\subseteq T'$, define
\[
 J_{T,T'}:\mathcal B_T\longrightarrow\mathcal B_{T'}
\]
to be the identity on objects and to send each arrow to the unique arrow
with the same endpoints. The theorem inclusion makes this functor
well-defined. It preserves the displayed Boolean operations, and
$J_{T',T''}J_{T,T'}=J_{T,T''}$. Consistency means that
$\mathcal B_T(\top,\bot)$ is empty; completeness means that at least one
of $\mathcal B_T(\top,\alpha)$ and
$\mathcal B_T(\top,\neg\alpha)$ is inhabited for each $\alpha$.

\begin{lemma}\label{num:effective-proofs}
Every c.e. theory $T$ extending $T_0$ admits a recursive rule
presentation extending $\Sigma_0$ and having theorem range $1_T$.
\end{lemma}
\begin{proof}
Fix an enumerator for $T$. Add to $\Sigma_0$ a nullary instance
$(\alpha,m)$ whenever $\alpha$ is printed within its first $m$ steps,
with conclusion $\alpha$. This instance condition is recursive.
Every added axiom belongs to $T$, and every member of $T$ occurs as an
axiom. Since $T$ is deductively closed, the new presentation has exactly
$T$ as its theorem set. Theorem~\ref{poly:initial} supplies its recursive
proof object. Inclusion of rule instances is a presentation morphism,
so the proof and range functors give $t^+\le1_T$.
\end{proof}

\begin{theorem}[Numerical essential incompleteness]\label{num:essential}
For the universal evaluator $\tau$ above, the theory $T_0$ of
\eqref{num:axioms} is consistent and c.e. Every consistent c.e. extension
$T$ of $T_0$ leaves some atom $P_n$ undecided, and is therefore incomplete.
\end{theorem}
\begin{proof}
The proof object $D$ is recursive and its conclusion map is computable,
so $T_0$ is c.e.; Lemma~\ref{num:consistent} gives consistency.
For an extension $T$, write $t_T=1_T$.
Lemma~\ref{num:effective-proofs} and monotonicity of inverse image give
\[
 t^+\le t_T,\qquad t^-\le\nu^*t_T.
\]
Consistency gives $t_T\nu^*t_T=0_S$. If $T$ were complete, these domains
would be complementary on $S$ and would separate $t^+,t^-$, contrary to
Proposition~\ref{num:proof-inseparable}.

More precisely, the domains
$\theta^*t_T$ and $\theta^*\nu^*t_T$ on $\N$ are disjoint and contain
$a,b$, respectively. If every $P_n$ were decided, their union would be
$1_{\N}$, contradicting inseparability of $a,b$.
\end{proof}

The valuation in Lemma~\ref{num:consistent} establishes consistency of
$T_0$. The extension argument uses consistency of $T$; it does not require
that the same valuation satisfy its additional theorems.

\subsection{Hoare interpretation}

Define $k:\N\rightharpoonup\N$ by simulating $d(n)$ and returning its
output if that output is $0$ or $1$, and diverging otherwise. Then
$k^{-1}(0)=A$ and $k^{-1}(1)=B$. Put $p_n=1_{\{n\}}$ and
$q_i=1_{\{i\}}$. Under the valuation used above,
\begin{align}
 P_n\text{ is true}
 &\quad\Longleftrightarrow\quad H_k(p_n,q_0)\ne\varnothing,
 \label{num:hoare-positive}\\
 \neg P_n\text{ is true}
 &\quad\Longleftrightarrow\quad K_k(p_n,q_1)\ne\varnothing.
 \label{num:hoare-negative}
\end{align}
The first assertion requires termination with output $0$. The second
holds when $k(n)$ is undefined or has output $1$, exactly the cases in
which $n\notin A$. The negative axioms arise from the latter output
case, which has finite computation evidence.

The second equivalence uses singleton preconditions. For example, if
$f(0)=0$ and $f(1)=1$, then both
$H_f(1_{\{0,1\}},q_0)$ and $K_f(1_{\{0,1\}},q_1)$ are empty.
Thus the complement of a total-correctness assertion on an arbitrary
precondition is not generally the corresponding partial-correctness
assertion.

\subsection{Arithmetic representation of recursive sets}

The representation hypothesis of
Theorem~\ref{ct:representation-inseparable} gives another application.
Fix an effective first-order language of arithmetic, an effective
classical deductive calculus, and a recursive coding of finite terms and
formulas. The sentence codes form a recursive set; let $s$ be its
partial identity. Define $\nu$ on that set by
$\nu(\ulcorner\alpha\urcorner)=\ulcorner\neg\alpha\urcorner$.

Choose a total recursive enumeration $u\mapsto\varphi_u(v)$ of formulas
whose only possible free variable is $v$, covering every such formula.
For instance, use the formula with code $u$ when it has this property,
and otherwise use $0=0$. Write $\underline n$ for the numeral denoting
$n$. Numeral substitution defines a total recursive map
\[
 \sigma:\N^2\longrightarrow\N,\qquad
 \sigma(u,n)=\ulcorner\varphi_u(\underline n)\urcorner.
\]
Recursion on $n$ constructs the numeral from zero and successor.
Recursion on the formula tree replaces its free occurrences of $v$ by
that closed term. Closedness prevents variable capture. The algorithm
therefore terminates jointly in $u,n$, and $s\sigma=\sigma$.

An arithmetical theory $B$ \emph{defines $C\subseteq\N$ two-sidedly
numeralwise} if some formula $\varphi(v)$ satisfies
\[
 n\in C\Longrightarrow B\vdash\varphi(\underline n),\qquad
 n\notin C\Longrightarrow B\vdash\neg\varphi(\underline n).
\]

\begin{corollary}\label{num:recursive-representation}
Let $B$ be a consistent arithmetical theory in the effective language
above. If $B$ defines every recursive subset of $\N$ two-sidedly
numeralwise, then every consistent c.e. extension of $B$ is incomplete.
\end{corollary}
\begin{proof}
For a consistent c.e. extension $T$, let $t$ be its theorem domain and
put $t^-=\nu^*t$. These domains are disjoint. A complementary pair of
c.e. domains on $\N$ consists of a recursive set $C$ and its complement.
Choose an index $u$ of a formula defining $C$ two-sidedly numeralwise
in $B$, and let $g:\N\to\N$ be the total constant map with value $u$.
The corresponding instances remain provable in $T$. Consistency excludes
the opposite instances, whence
\[
 [\sigma(g\boxtimes1)]^*t=\pi_2^*1_C,
 \qquad
 [\sigma(g\boxtimes1)]^*t^-=\pi_2^*1_{\N\setminus C}.
\]
These are the representation equations of
Theorem~\ref{ct:representation-inseparable}. Hence $t,t^-$ are
inseparable relative to $s$, and cannot themselves be complementary.
\end{proof}

\subsection{Effective interpretations}

\begin{proposition}\label{num:interpretation}
Let $A,B$ be the inseparable pair above, and let $\mathcal T$ be a
consistent c.e. classical theory in an effective language. Suppose there
is a computable sequence of sentences $(\rho_n)_{n\in\N}$ such that
\[
 n\in A\Longrightarrow\mathcal T\vdash\rho_n,\qquad
 n\in B\Longrightarrow\mathcal T\vdash\neg\rho_n.
\]
Every consistent c.e. extension of $\mathcal T$ leaves some $\rho_n$
undecided.
\end{proposition}
\begin{proof}
The assignment $P_n\mapsto\rho_n$ extends by the Boolean constructors
to a computable translation $I$ of formulas. For a consistent c.e.
extension $\mathcal T'$, the set
\[
 T_I=\{\alpha\in S:\mathcal T'\vdash I(\alpha)\}
\]
is c.e., contains $\Gamma_0$, and is closed under tautologies and modus
ponens. It is consistent because $I$ preserves falsity. Apply
Theorem~\ref{num:essential} to $T_I$.
\end{proof}

An interpretation by computation formulas must therefore establish both
families of provable instances. The computability of the sentence map
includes effective coding and substitution.

For semantic predicate transformation, fix an arithmetic postcondition
$Q(y)$. Let $\operatorname{Trace}_f(n,c,y)$ express that $c$ is a halting
computation of $f$ from $n$ with output $y$. After effective
renaming of bound variables, the formulas
\[
 \exists c\,\exists y\,
   (\operatorname{Trace}_f(n,c,y)\land Q(y)),
 \qquad
 \forall c\,\forall y\,
   (\operatorname{Trace}_f(n,c,y)\Rightarrow Q(y))
\]
define $\operatorname{wp}_f(Q)$ and $\operatorname{wlp}_f(Q)$ in the
standard natural numbers. Equation~\eqref{hp:halting-predicates} shows
why constructing these formulas cannot decide their truth.
Cook's effective-verification obstruction uses the partial-correctness
assertion $\{\mathrm{true}\}\,C\,\{\mathrm{false}\}$, which holds
exactly when $C$ terminates on no initial state
\cite[\S6, p.~85]{Cook}. His relative completeness theorem assumes an
expressive assertion language and a semantically complete assertion
proof system \cite[\S6, p.~86]{Cook}.

\section{Logical presentations in recursion categories}\label{prob:section}

\begin{problem}\label{prob:intrinsic}
Let $\R$ be a recursion category with ranges and with the binary case
distinction required by Theorem~\ref{tur:inseparable}. Construct a
category of logical presentations in a specified completion of $\R$,
with assertion objects, negation maps, polynomial functors for inference,
and initial algebras of finite proofs. Construct, from the computational
inseparable pair, the positive and negative proof maps of
Theorem~\ref{ct:certificate-transfer}. Establish a consistent base
presentation and identify the morphisms that represent effective
extensions.
\end{problem}

The proof of Theorem~\ref{poly:initial} uses recursive
finite-tree coding and a terminating fold into every algebra. An abstract
analogue needs corresponding existence and recursion principles in the
chosen completion.

Niu and Spivak ask whether propositions in a functor category associated
with polynomial trees admit a G\"odel-like coding into the language of
the corresponding dynamical systems
\cite[Chapter~9, Question~14, p.~350]{NS2024}.
For the existence problem above, the relevant universal property is
the initial-algebra property of Theorem~\ref{poly:initial}.

Logical negation in a presentation need not complement a computational
domain. The effective map $\nu$ acts on assertion codes, and inverse
image along $\nu$ selects the assertions whose negations are provable.
If a completion adds Boolean complements, computational separation must
still refer to the original domains: in a Boolean algebra, every
disjoint pair is separated by a complemented element.

An interpretation in Robinson arithmetic or another specified arithmetic
theory requires a computable translation of assertions and proofs with
the properties stated in Proposition~\ref{num:interpretation}.
Corollary~\ref{num:recursive-representation} gives a separate sufficient
condition in terms of numeralwise definability of recursive sets.

\begin{thebibliography}{99}
\bibitem{NS2024} N.~Niu and D.~I.~Spivak,
\emph{Polynomial Functors: A Mathematical Theory of Interaction},
draft of 17 July 2024.
\url{https://github.com/ToposInstitute/poly}.
\bibitem{DPH} R.~A.~Di Paola and A.~Heller,
Dominical categories: recursion theory without elements,
\emph{Journal of Symbolic Logic} \textbf{52}(3) (1987), 594--635.
\href{https://doi.org/10.2307/2274352}{doi:10.2307/2274352}.
\bibitem{Heller} A.~Heller,
An existence theorem for recursion categories,
\emph{Journal of Symbolic Logic} \textbf{55}(3) (1990), 1252--1268.
\href{https://doi.org/10.2307/2274486}{doi:10.2307/2274486}.
\bibitem{CockettLack} J.~R.~B.~Cockett and S.~Lack,
Restriction categories I: categories of partial maps,
\emph{Theoretical Computer Science} \textbf{270}(1--2) (2002), 223--259.
\bibitem{CockettGarner} J.~R.~B.~Cockett and R.~Garner,
Restriction categories as enriched categories,
\emph{Theoretical Computer Science} \textbf{523} (2014), 37--55.
\href{https://arxiv.org/abs/1211.6170}{arXiv:1211.6170}.
\bibitem{Dijkstra} E.~W.~Dijkstra,
Guarded commands, non-determinacy and a calculus for the derivation of programs,
EWD418, author manuscript.
\url{https://www.cs.utexas.edu/~EWD/transcriptions/EWD04xx/EWD418.html}.
\bibitem{DijkstraScholten} E.~W.~Dijkstra and C.~S.~Scholten,
\emph{Predicate Calculus and Program Semantics},
Springer, New York, 1990.
\href{https://doi.org/10.1007/978-1-4612-3228-5}{doi:10.1007/978-1-4612-3228-5}.
\bibitem{Cook} S.~A.~Cook,
Soundness and completeness of an axiom system for program verification,
\emph{SIAM Journal on Computing} \textbf{7}(1) (1978), 70--90.
\url{https://www.cs.toronto.edu/~sacook/homepage/soundness.pdf}.
\bibitem{ArthanEtAl} R.~Arthan, U.~Martin, E.~A.~Mathiesen, and P.~Oliva,
A general framework for sound and complete Floyd--Hoare logics,
\emph{ACM Transactions on Computational Logic} \textbf{11}(1) (2009),
Article~7, 1--31.
\href{https://doi.org/10.1145/1614431.1614438}{doi:10.1145/1614431.1614438};
\href{https://arxiv.org/abs/0807.1016}{arXiv:0807.1016}.
\end{thebibliography}
\end{document}
